\documentclass[aps,prd,onecolumn,nofootinbib,superscriptaddress]{revtex4-2}

% Font and encoding settings
\usepackage[T2A]{fontenc}
\usepackage[utf8]{inputenc}
\usepackage[english,ukrainian]{babel} % Українська — основна

% Mathematical packages
\usepackage{amsmath,amssymb,amsfonts}
\usepackage{graphicx}
\usepackage{tikz}
\usepackage{float}
\usepackage{csquotes}
\usepackage{enumitem}

% hyperref is recommended to be loaded last
\usepackage[unicode]{hyperref}

% ----------------------------------------------------------------
% Metadata and reusable commands
% ----------------------------------------------------------------
\newcommand{\papertitleuk}{Поле відмінностей та емерджентні інваріанти: концептуальна нотатка про суб’єкт}
\newcommand{\papertitleen}{The Difference Field and Emergent Invariants: A Conceptual Note on the Subject}
\newcommand{\paperauthor}{Віталій Дурнєв}
\newcommand{\paperauthorlatin}{Vitalii Durniev}

\hypersetup{
    pdftitle={\papertitleuk},
    pdfauthor={\paperauthor},
    pdfsubject={Концептуальна рамка виникнення простору, часу, інформації, інваріантів і суб’єкта з поля відмінностей},
    pdfkeywords={поле відмінностей, інваріанти, згортання, суб’єкт, суб’єктність, простір, час, інформація, кібернетика, філософія свідомості},
    pdfcreator={LaTeX with hyperref}
}

\begin{document}

% У REVTeX 4.2 дані про статтю описуються після \begin{document}
% ----------------------------------------------------------------
% Article metadata
% ----------------------------------------------------------------
\title{%
\papertitleuk\\[0.5em]
\large\itshape \papertitleen
}

\thanks{Ця робота поширюється на умовах ліцензії CC BY-NC-ND 4.0. Комерційне використання та поширення змінених версій без письмового дозволу автора заборонено.}

\author{\paperauthor}
\affiliation{Незалежний дослідник, Київ, Україна}

\date{17 квітня 2026}

\begin{abstract}
У статті пропонується концептуальна рамка, в якій світ, доступний суб’єкту (S), не ототожнюється з повнотою реальності, а розглядається як система інваріантних структур ($\mathcal{I}$), доступних для сприйняття та узгодження. Вихідною точкою слугує граничний уявний експеримент, у якому суб’єкт (S), поміщений у майже однорідне середовище, фізично насичене даними, але позбавлене доступних зовнішніх опор для стабільного розрізнення. На цій основі вводиться поняття поля відмінностей (D) як дооб’єктної основи, що не збігається ані з готовим світом об’єктів, ані з інформацією для конкретного суб’єкта (S).

Показано, що інформація для суб’єкта (S) виникає лише там, де дані стають доступними у формі інваріантних структур ($\mathcal{I}$), а простір і час можуть бути витлумачені як похідні режими узгодження таких структур, а не як безумовно задані передумови досвіду. Далі окреслюється гіпотеза багаторівневого згортання, у межах якої суб’єкт (S) розглядається як інваріантна структура вищого порядку, що не пізнає поле відмінностей (D) у повноті, а узгоджує доступну частину поля, зберігаючи власну цілісність.

У цьому сенсі поняття суб’єкта \(S\) не зводиться до простої точки спостереження, а наближається до ідеї цілісної моделі самості, хоча й у межах іншої концептуальної мови~\cite{Metzinger2003}.

Стаття має характер першої концептуальної презентації запропонованої ідеї й розглядається як запрошення до критичної міждисциплінарної дискусії.
\end{abstract}

\maketitle

\section{Вступ}
Звичний науковий опис світу, як правило, виходить із того, що існують об’єкти, розміщені в просторі й часі, а суб’єкт (S) лише фіксує або дедалі точніше описує вже наявну структуру. Однак така картина може бути не первинною, а похідною. Те, що здається нам готовим світом речей, подій, відстаней і тривалостей, може бути результатом певного способу стабілізації відмінностей, доступного конкретному типу суб’єкта (S).

Ця стаття виходить із припущення, що для побудови фундаментальнішої рамки слід розрізнити щонайменше три рівні: те, що існує незалежно від конкретного суб’єкта (S); те, що може бути стійко виділене у формі інваріантів (I) або систем інваріантів ($\mathcal{I}$); і те, що стає узгодженим світом для суб’єкта (S). Таке розрізнення важливе, оскільки в повсякденному та науковому мисленні ці рівні зазвичай неусвідомлено зливаються. Ми схильні ототожнювати реальність із доступним нам світом, а світ — із набором уже оформлених об’єктів. Проте немає підстав вважати, що будь-який можливий суб’єкт (S) мав би бачити той самий світ у тій самій структурі.

Щоб зробити цю проблему максимально наочною, у статті розглядається граничний уявний експеримент, у якому суб’єкт (S), якого помістили у майже однорідне середовище, де фізично присутня велика кількість даних, але майже відсутні доступні зовнішні інваріанти (I). Уявний характер цього прикладу частково натхненний перцептивною ситуацією, близькою до тієї, яку створює інсталяція Antony Gormley *Blind Light* (2007). Мета цього експерименту не в тому, щоб емпірично довести нову теорію, а в тому, щоб виявити межу інтуїції, за якою світ автоматично мислиться як уже структурований.

Головна теза роботи полягає в тому, що наявність даних ще не гарантує наявності світу для суб’єкта (S). Світ виникає лише там, де дані можуть бути стійко виділені, зіставлені й інтегровані в узгоджену систему інваріантів ($\mathcal{I}$). На цій основі в статті вводиться поняття поля відмінностей (D), яке розглядається як дооб’єктна основа, що не збігається ні з інформацією для конкретного суб’єкта (S), ні з уже побудованими простором і часом. Далі висувається гіпотеза, що суб’єкт (S) може бути витлумачений як інваріант вищого порядку (I), який не охоплює поле (D) у повноті, а узгоджує його доступну частину, зберігаючи можливість переживати себе як цілісну, хоча й інваріантну, структуру.

Стаття не претендує на статус завершеної фізичної або математичної теорії. Її мета скромніша: запропонувати концептуальну мову, яка дозволяє спільно розглядати поле відмінностей D, інваріантні структури I, багаторівневе згортання і суб’єкта S. Методологічно запропонована рамка близька до системної традиції, в якій реальність не редукується до набору окремих речей, а розглядається як поле відношень, відмінностей і рівнів організації, з яких постають стійкі форми світу~\cite{Bateson1972}. У цій перспективі суб’єкт S може бути витлумачений як інваріантна структура вищого порядку, що підтримує узгоджений доступний світ через відбір, згортання та стабілізацію інваріантних структур I. Така інтерпретація, своєю чергою, частково перегукується із сучасними моделями, у яких когнітивна система зберігає свою цілісність шляхом мінімізації невизначеності~\cite{Friston2010}.
\section{Граничний експеримент: кімната з однорідним туманом}
Уявімо кімнату, повністю заповнену густим, сухим і однорідно освітленим туманом. Світло в ній рівномірне: не існує жодної видимої межі, джерела світла, напрямку чи контрасту. Суб’єкт S перебуває в центрі цієї кімнати, не торкається стін і не має змоги побачити власне тіло. Його руки опущені, а туман настільки густий, що суб’єкт S не фіксує навіть наявності власного носа у поточному перцептивному досвіді. Середовище при цьому настільки однорідне, що будь-які локальні особливості зникають у загальному фоні.

Фізично така кімната зовсім не є порожньою. У ній присутня велика кількість молекул води, фотонів, локальних неоднорідностей і можливих конфігурацій. Теоретично з певних молекул можна було б виділити коло, з інших — інше коло меншого розміру; їхнє порівняння могло б створити відношення, подібне до просторового. Не змінюючи положення молекул, а лише вибираючи їх у певному порядку, можна було б уявно відтворити довільно складні структури — наприклад, повний текст Encyclopaedia Britannica. Однак для суб’єкта (S), який перебуває в кімнаті, ці можливості залишаються недоступними. Він не розрізняє окремих молекул, не бачить їхніх взаємних відношень і не має способу стабільно виділити будь-яку із цих потенційних структур.

У результаті виникає особливий досвід: навколо є фізично насичене середовище, але майже немає світу як системи доступних об’єктів і відношень. Суб’єкт (S) не може надійно виділити жодної зовнішньої точки, до якої можна було б прив’язати систему координат. Він не має достатньо підстав, щоб побудувати навіть найпростішу геометрію. Якщо припустити також відсутність сталої гравітаційної орієнтації, наприклад у режимі рівномірного падіння, то зникає й останній зовнішній напрямок, який можна було б використати для відліку.

Попри це, повна відсутність досвіду не настає. У суб’єкта (S) зберігаються внутрішні ритми: серцебиття, дихання, залишкова тілесна чутливість. Саме вони стають майже єдиними доступними інваріантними структурами ($\mathcal{I}$). Через них зберігається мінімальне переживання власної безперервності, навіть якщо зовнішній світ майже зникає.

Особливо важливим є те, що в такій ситуації стає проблематичним саме поняття часу. Якщо зовнішні події не виділяються, а зміни не можуть бути стабільно співвіднесені, то суб’єкт (S) не має з чого побудувати послідовність. Без доступних циклів і повторюваних відмінностей зміни не стають подіями, а без подій не виникає й час у звичному сенсі. Єдиною основою для переживання тривалості лишаються внутрішні ритми організму. Це має очевидний зв’язок із підходами embodied cognition~\cite{VarelaThompsonRosch1991}.

Аналогічно, простір також перестає бути чимось безумовно даним. Простір у звичному розумінні потребує хоча б двох стабільно співвіднесених інваріантних структур ($\mathcal{I}$), між якими можна встановити відношення відстані, напрямку або межі. У кімнаті з однорідним туманом таких структур майже немає. Тому простір не зникає як фізична можливість, але перестає існувати як доступна форма організації досвіду для конкретного суб’єкта (S) який перебуває в цій кімнаті.

Цей thought experiment дозволяє зробити перший принциповий висновок: наявність фізичних даних ще не гарантує наявності світу для суб’єкта (S). Світ виникає лише там, де існують інваріантні структури ($\mathcal{I}$), які можуть бути стабільно виділені, повторно співвіднесені й інтегровані в узгоджену картину для суб’єкта (S). Саме ця відмінність між фізичною насиченістю середовища та відсутністю доступного світу і є відправною точкою для введення поняття поля відмінностей (D).

\section{Дані, інформація і доступний світ}
Приклад із кімнатою показує, що між фізичною наявністю середовища і світом, який переживає суб’єкт (S), існує принципова різниця. Щоб зробити цю різницю явною, доцільно розрізнити щонайменше три поняття: дані, інформацію і доступний світ.

Під даними в цій статті розумітиметься вся сукупність відмінностей, конфігурацій і потенційних співвідношень, що фізично присутні незалежно від того, чи здатен конкретний суб’єкт (S) їх виділити. У кімнаті з туманом даних надзвичайно багато: молекули, фотони, локальні неоднорідності, можливі геометричні конфігурації. Проте сама по собі наявність даних ще не означає, що вони мають значення для конретного суб’єкта (S).

Інформація виникає лише тоді, коли деякі дані можуть бути виділені як інваріантні структури ($\mathcal{I}$), тобто як форми, які зберігають достатню стійкість для повторного розрізнення і співвіднесення. Якщо суб’єкт (S) не має способу стабільно виділити певну структуру, то для нього вона не є інформацією, навіть якщо вона фізично присутня. У кімнаті з туманом суб’єкт (S) теоретично оточений величезною кількістю можливих структур, але майже жодна з них не стає інформацією, оскільки не може бути повторно впізнана суб’єктом (S).

\begin{quote}
\textbf{Визначення.}
Інформація для суб’єкта \(S\) — це дані, які стали доступними у формі інваріантних структур \(\mathcal{I}\), придатних для повторного співвіднесення.
\end{quote}

Таке визначення істотно відрізняється від підходів, у яких інформація розглядається як об’єктивна властивість будь-якої фізичної конфігурації. У межах запропонованої рамки одна й та сама конфігурація може бути інформацією для одного типу суб’єкта (S) і не бути нею для іншого. Інформація не є просто “тим, що є”; вона є тим, що може бути інваріантно доступним для конкретного суб’єкта (S).

Доступний світ, у свою чергу, є не всією множиною інформації, а узгодженою системою інваріантних структур ($\mathcal{I}$), які можуть бути інтегровані в стійку картину. Для людини в звичайних умовах доступний світ включає просторові об’єкти, межі, рух, причинність, часову послідовність. У кімнаті з однорідним туманом цей світ майже руйнується: більшість звичних зовнішніх інваріантних структур ($\mathcal{I}$) стає недоступною, і тому зовнішній світ перестає бути організованим.

Проте навіть у цій ситуації доступний світ не зникає повністю. Внутрішні ритми організму — серцебиття, дихання, зміна напруги тіла — залишаються доступними інваріантними структурами ($\mathcal{I}$). Саме вони утворюють мінімальний залишок світу для суб’єкта (S). Отже, доступний світ не є тотожним зовнішньому середовищу; він визначається тим, які інваріантні структури ($\mathcal{I}$) суб’єкт (S) здатен утримувати та інтегрувати.

Таким чином, між даними, інформацією і світом існує ієрархія:

\[
\text{дані} \supset \text{інформація для } S \supset \text{доступний світ}
\]

Дані охоплюють усе фізично можливе. Інформація становить лише ту частину даних, яка стала доступною як інваріантна структура. Доступний світ є ще вужчим: це не просто множина окремих інваріантів (I), а їхня узгоджена система ($\mathcal{I}$).

Саме ця ієрархія пояснює, чому суб’єкт (S) у кімнаті переживає щось близьке до відсутності світу, хоча фізично він перебуває в середовищі, надзвичайно багатому на дані. Те, що переживається як “нічого”, є не відсутністю фізичної реальності, а відсутністю доступних інваріантних структур ($\mathcal{I}$), достатніх для побудови світу. Наступний крок полягає в тому, щоб ввести поняття поля відмінностей (D), яке дозволяє мислити дані як більш фундаментальний рівень, ніж будь-який уже сформований світ.

\section{Поле відмінностей \texorpdfstring{$D$}{D}}
Якщо поле (D) саме по собі ще не є світом, то постає питання: що саме повинно виникнути, щоб світ став можливим для суб’єкта (S)? У межах запропонованої рамки мінімальною умовою цього є поява інваріантних структур ($\mathcal{I}$).

Під інваріантною структурою в цій статті розумітиметься така форма, яка може бути повторно виділена і співвіднесена, незважаючи на локальні зміни або варіації. Інваріантна структура не обов’язково є абсолютно незмінною. Навпаки, вона може складатися з багатьох відмінних конфігурацій, які на більш високому рівні згортання сприймаються як “те саме”.

Наприклад, те, що людина називає одним і тим самим ритмом серця, ніколи не повторюється абсолютно точно. Кожен удар серця відрізняється від попереднього. Проте ця послідовність може бути згортана в одну інваріантну структуру (I) — “серцебиття”. Аналогічно, в повсякденному досвіді суб’єкт (S) сприймає “той самий” предмет навіть тоді, коли освітлення, положення, кут огляду або дрібні деталі змінюються.

Отже, інваріантна структура (I) не є тотожністю в строгому сенсі. Вона є результатом стискання множини відмінностей у форму, достатньо стійку для повторного співвіднесення.

\begin{quote}
\textbf{Визначення.}
Інваріантна структура (I) — це згортання множини відмінностей поля (D) у форму, яка може бути повторно впізнана і співвіднесена.
\end{quote}


Саме інваріантні структури ($\mathcal{I}$) роблять можливим виникнення світу. Без них суб’єкт (S) не може відрізнити “те саме” від “іншого”, не може побудувати межу, об’єкт або послідовність. У кімнаті з однорідним туманом зовнішній світ майже зникає саме тому, що середовище не дає достатньо доступних інваріантних структур ($\mathcal{I}$). Суб’єкт (S) не може стабільно виділити ні напрямок, ні межу, ні форму.

Попри це, деякі інваріантні структури ($\mathcal{I}$) все ж залишаються. Передусім це внутрішні ритми організму: серцебиття, дихання, тілесна напруга. Вони є мінімальними структурами, які суб’єкт (S) може повторно співвідносити. Саме тому навіть у ситуації майже повної втрати зовнішнього світу не зникає повністю переживання власної безперервності.

У межах цієї статті доцільно розрізняти два типи інваріантних структур:

\begin{enumerate}
    \item \textbf{Зовнішні інваріантні структури} --- ті, що пов’язані з відносно стабільними особливостями середовища;

    \item \textbf{Внутрішні інваріантні структури} --- ті, що пов’язані з самим суб’єктом \(S\) і підтримують його мінімальну цілісність.
\end{enumerate}

У звичайному людському досвіді ці два типи тісно переплетені. Зовнішній світ постійно співвідноситься з внутрішніми ритмами та станами, а внутрішня цілісність підтримується через взаємодію із зовнішніми інваріантними структурами ($\mathcal{I}$). У кімнаті з туманом цей баланс радикально зміщується: майже всі зовнішні інваріантні структури зникають, і тому суб’єкт (S) залишається майже сам на сам із власними внутрішніми інваріантами.

Особливо важливо, що інваріантні структури ($\mathcal{I}$) не лише дозволяють суб’єкту (S) побачити світ, а й визначають, який саме світ для нього можливий. Інший тип суб’єкта (S), який мав би доступ до інших інваріантних структур ($\mathcal{I}$), міг би будувати зовсім іншу картину реальності. Те, що для людини є об’єктом, простором або подією, може не існувати для іншого типу суб’єкта (S), і навпаки.

Звідси випливає, що світ не є простою сукупністю речей, однаково доступних для всіх можливих суб’єктів (S). Світ є системою інваріантних структур ($\mathcal{I}$), доступних конкретному типу суб’єкта (S).

Саме тому інваріантні структури ($\mathcal{I}$) не слід розглядати як додаток до вже готового світу. Вони є умовою його виникнення. Без них існує лише поле (D), насичене відмінностями, але ще не організоване у світ.

\section{Інваріантні структури як умова світу}
Якщо дані не збігаються з інформацією для суб’єкта (S), а інформація не збігається з доступним світом, то постає питання: що саме є тим рівнем, на якому дані існують до будь-якого виділення, інтерпретації чи узгодження? У цій статті пропонується позначати цей рівень як поле відмінностей (D).

\begin{quote}
\textbf{Визначення.}
Поле відмінностей \(D\) — це дооб’єктна сукупність усіх можливих відмінностей і співвідношень, які ще не організовані у світ об’єктів, простору, часу або смислів.
\end{quote}

У цьому сенсі поле (D) є дооб’єктною основою. Воно не містить простір і час як уже задані форми. Простір і час можуть виникати лише пізніше, як результат узгодження інваріантних структур ($\mathcal{I}$). Те, що в звичайному досвіді виглядає як тривимірний простір або послідовність подій, у межах запропонованої рамки не є первинною властивістю (D), а лише одним із можливих способів його організації.

Важливо також, що (D) не є пустотою. Навпаки, кімната з однорідним туманом показує, що можлива ситуація, у якій суб’єкт (S) переживає майже повну відсутність світу, хоча навколо залишається надлишок фізичних даних. Тому переживання “нічого” не слід тлумачити як відсутність поля (D). Це лише означає, що для даного суб’єкта (S) в даних умовах майже не існує доступних інваріантних структур ($\mathcal{I}$).

Поле (D) не повинно також мислитися як однорідна маса без внутрішньої організації. Навіть у кімнаті з туманом теоретично можливо виділити з молекул певні конфігурації: наприклад, коло або іншу форму. Але ці конфігурації не є вже готовими інваріантними структурами. Вони існують лише як потенційна можливість. Саме тому (D) можна описати як поле не готових об’єктів, а потенційної розрізненості.

У межах цієї статті пропонується виходити з припущення, що в полі (D) не існує двох абсолютно тотожних елементів або конфігурацій. Тотожність виникає лише на рівні подальшого узгодження, коли деякі відмінності починають стискатися в одну інваріантну структуру (I). Те, що для суб’єкта (S) виглядає як “те саме”, на рівні (D) може залишатися сукупністю відмінних, але достатньо близьких конфігурацій.

Таким чином, однаковість не є фундаментальною властивістю реальності, а результатом згортання відмінностей. Саме завдяки цьому пізніше стають можливими повторення, об’єкти і стійкі світи.

Поле (D) не є світом для жодного конкретного суб’єкта (S). Жоден суб’єкт (S) не може повністю охопити (D) або утримати всі його відмінності в межах власного досвіду. Причина цього полягає не лише в обмеженості суб’єкта, а й у самій природі (S): щоб бути цілісною інваріантною структурою, суб’єкт (S) змушений стискати, відкидати і впорядковувати більшу частину (D). Отже, будь-який світ, доступний суб’єкту (S), є лише локальним і частковим проявом поля відмінностей.

Це дозволяє сформулювати центральну тезу:
\begin{quote}
(D) не є світом; світ є інваріантно доступним проявом (D) для конкретного суб’єкта (S).
\end{quote}

Наступним кроком є пояснення того, як із поля (D) можуть виникати інваріантні структури ($\mathcal{I}$) і чому саме вони стають мінімальною умовою появи світу для суб’єкта (S).

\section{Простір і час як похідні режими узгодження}
У звичному досвіді простір і час здаються найфундаментальнішими характеристиками світу. Ми схильні припускати, що об’єкти завжди вже розташовані в просторі, а події завжди вже відбуваються в часі. Проте приклад із кімнатою показує, що ця інтуїція може бути похідною.

У кімнаті з однорідним туманом фізично присутнє середовище, але суб’єкт (S) не має достатньо доступних інваріантних структур ($\mathcal{I}$), щоб побудувати навіть найпростішу геометрію. Простір у звичному сенсі потребує хоча б двох стабільно співвіднесених інваріантних структур (I), між якими можна встановити відношення: ближче чи далі, вище чи нижче, ліворуч чи праворуч. Якщо таких структур немає, то немає й підстав для виникнення координат, меж або напрямків.

Отже, простір не є просто “контейнером”, який існує незалежно від будь-якого способу сприйняття. Він виникає лише там, де інваріантні структури ($\mathcal{I}$) можуть бути достатньо стійко співвіднесені одна з одною. Простір у цьому сенсі є способом організації доступних відмінностей, а не первинною властивістю поля (D).

Аналогічно, час також не може бути виведений безпосередньо з самого факту зміни. Щоб виник час, недостатньо, щоб щось відбувалося. Потрібно, щоб різні стани могли бути співвіднесені як повторення, послідовність або цикл. Якщо суб’єкт (S) не має доступу до інваріантних структур ($\mathcal{I}$), які можна зіставляти, то зміни не стають подіями, а без подій не виникає й час у звичному сенсі.

У кімнаті з туманом це особливо помітно. Зовнішні зміни, навіть якщо вони існують, майже не можуть бути інтегровані в послідовність. Тому єдиним джерелом переживання тривалості стають внутрішні інваріантні структури ($\mathcal{I}$): серцебиття, дихання, зміни стану тіла. Саме вони створюють мінімальну основу для того, що суб’єкт (S) може переживати як “до” і “після”.

Звідси випливає, що час не є незалежною субстанцією або фоном, у якому відбуваються події. Час виникає як результат узгодження інваріантних структур ($\mathcal{I}$), які можуть бути співвіднесені в послідовність.

У межах запропонованої рамки можна сформулювати такі визначення:

\begin{quote}
\textbf{Визначення.}
Простір — це спосіб узгодження співіснуючих інваріантних структур ($\mathcal{I}$).
\end{quote}

\begin{quote}
\textbf{Визначення.}
Час — це спосіб узгодження послідовних інваріантних структур ($\mathcal{I}$).
\end{quote}

Обидва ці режими не існують на рівні поля (D) як такого. Вони виникають лише тоді, коли деяка частина (D) стає доступною у формі інваріантних структур ($\mathcal{I}$) і може бути інтегрована в стійку картину.

Тому простір і час не є первинними передумовами світу. Навпаки, вони самі є похідними від більш фундаментального процесу — появи, співвіднесення і згортання інваріантних структур ($\mathcal{I}$). Саме тому різні типи суб’єкта (S), які мають доступ до різних інваріантних структур ($\mathcal{I}$), потенційно могли б будувати зовсім інші форми простору і часу або не будувати їх зовсім.

Приклад із кімнатою демонструє це у найпростішому вигляді: фізичне середовище залишається, але простір і час майже зникають як доступні форми організації досвіду. Це означає, що вони не можуть бути прийняті як безумовно дані властивості реальності. Вони є лише одним із можливих способів, у який певний тип суб’єкта (S) організує доступну йому частину поля (D).

\section{Багаторівневе згортання}
Якщо поле (D) не є вже готовим світом, а інваріантні структури ($\mathcal{I}$) є умовою виникнення світу для суб’єкта (S), то постає питання: яким чином такі структури можуть виникати з поля відмінностей? У межах цієї статті висувається гіпотеза, що між (D) і доступним світом існує послідовність рівнів згортання.

Нехай (D) є первинним полем відмінностей, у якому ще не існує ані простору, ані часу, ані готових інваріантних структур ($\mathcal{I}$). Тоді перше згортання породжує новий рівень, який позначимо як (D'). (D') не є просто підмножиною (D). Це нове поле, що виникає як результат стискання частини відмінностей (D) у форми, придатні для повторного співвіднесення.

Схематично це можна записати як:
\[C(D)=D'\]

де (C) позначає оператор згортання.

Оператор (C) не усуває відмінності повністю. Його дія полягає в тому, що множина різних конфігурацій починає розглядатися як одна інваріантна структура (I). Те, що на рівні (D) існує як безліч відмінних конфігурацій, на рівні (D') може бути згорнуте в одну форму, яка стає доступною для повторного співвіднесення.

Принципово важливо, що (D) при цьому не зникає. Згортання не є знищенням попереднього рівня, а лише появою нового рівня організації. Тому (D) і (D') співіснують.

Подальші акти згортання можуть породжувати нові рівні:

\[ D \to D' \to D'' \to D''' \to \dots \]

Кожен новий рівень характеризується меншою потужністю опису. Те, що на нижчому рівні виглядає як величезна кількість відмінностей, на вищому рівні може бути згорнуте в невелику кількість інваріантних структур ($\mathcal{I}$).

Ідея багаторівневого згортання, в якій кожен наступний рівень зменшує різноманітність, але підвищує стійкість і інтегрованість форм, частково перегукується з уявленнями про багатомасштабну організацію і самоподібність~\cite{Mandelbrot1982,Mandelbrot1977}.

Таким чином, кожне згортання одночасно:
\begin{itemize}
    \item зменшує кількість доступних відмінностей;
    \item підвищує стійкість інваріантних структур \(I\);
    \item робить можливим новий рівень узгодження.
\end{itemize}

Саме завдяки цьому стає можливою поява дедалі складніших інваріантних структур ($\mathcal{I}$). Рівень (D') може містити найпростіші структури повторення. Рівень (D'') — вже їхні стійкі поєднання. На ще вищих рівнях можуть виникати структури, які здатні не лише підтримувати власну цілісність, а й співвідносити себе з іншими структурами.

Важливо, що кожен новий рівень спирається лише на частину попереднього. Не всі відмінності (D) потрапляють у (D'), і не всі структури (D') переходять у (D''). Тому будь-яке згортання є одночасно актом відбору і втрати. Воно не відкриває весь можливий світ, а лише певний спосіб його організації.

Саме тут стає можливим пояснення того, чому різні типи суб’єкта (S) можуть бачити різні світи. Якщо певний тип суб’єкта (S) виникає на основі специфічної послідовності згортань, то саме ця послідовність визначає, які інваріантні структури ($\mathcal{I}$) для нього доступні, а які залишаються невидимими. Світ \emph{Homo sapiens} є лише одним із варіантів згортання.

Разом з тим різні послідовності згортань можуть частково перекриватися. Тому різні суб’єкти (S) можуть мати спільні елементи світу, навіть якщо їхні інваріантні структури ($\mathcal{I}$) не збігаються повністю. Саме це може пояснювати, чому люди, маючи різні індивідуальні досвіди, все ж бачать приблизно один і той самий світ.

Особливо важливо, що нижчі рівні не перестають впливати на вищі. Навіть якщо деяка структура вже стала інваріантною на рівні (D'') або (D'''), вона продовжує залежати від нижчих рівнів, на яких залишається значно більше відмінностей. Це означає, що між різними інваріантними структурами верхнього рівня можуть зберігатися зв’язки через спільну нижчу основу. У подальших роботах така багаторівнева взаємодія може стати основою для інтерпретації фізичних полів, нелокальних ефектів або дальнодії.

У межах цієї статті достатньо сформулювати головний принцип:

\begin{quote}
Світ, доступний суб’єкту (S), не є прямим образом поля (D), а результатом послідовного багаторівневого згортання, у якому кожен рівень зменшує різноманітність і підвищує інваріантність.
\end{quote}

Тепер поняття суб’єкта (S), введене інтуїтивно на початку статті, можна точніше описати як інваріантну структуру вищого порядку.

\section{Суб’єкт як узгоджена інваріантна структура}
У межах запропонованої рамки суб’єкт (S) не може бути розглянутий як щось первинне і самодостатнє. Якщо світ виникає лише там, де з поля (D) виділяються та узгоджуються інваріантні структури ($\mathcal{I}$), то саме (S) також повинен бути результатом такого згортання.

Це означає, що (S) не передує інваріантним структурам ($\mathcal{I}$), а виникає разом із ними. Водночас (S) не є просто однією з інваріантних структур серед інших.

З огляду на це (S) можна визначити як:
\begin{quote}
\textbf{Визначення.}
Суб’єкт (S) — це інваріантна структура вищого порядку (I), яка інтегрує доступні інваріантні структури ($\mathcal{I}$) в узгоджений світ і при цьому підтримує власну інваріантність внутрішньої моделі світу.
\end{quote}

Таке визначення означає, що (S) ніколи не охоплює все поле (D). Навпаки, він завжди виникає на основі лише певної послідовності згортань і лише певної множини інваріантних структур ($\mathcal{I}$). Те, що не потрапляє в цю послідовність, не зникає, але не стає частиною доступного світу для даного суб’єкта (S).

Отже, будь-який суб’єкт (S) є локальним і частковим. Він не пізнає (D) у повноті, а лише вибудовує стійку модель тієї частини (D), яка виявилася для нього доступною. Саме тому різні суб’єкти (S) можуть мати різні світи, навіть якщо вони виникають на основі спільного поля відмінностей.

Разом з тим суб’єкти (S) роду \emph{Homo sapiens} зазвичай сприймають світ подібно. У межах цієї статті це можна пояснити тим, що (S) роду \emph{Homo sapiens} виникають на основі близьких або частково перекривних послідовностей згортань. Вони мають подібні тіла, подібні внутрішні ритми, подібні сенсорні обмеження, і тому доступні їм інваріантні структури ($\mathcal{I}$) значною мірою збігаються. Саме завдяки цьому люди можуть говорити про “один і той самий” світ, хоча жоден суб’єкт (S) не має до нього повного доступу і абсолютна тотожність світів для різних суб’єктів є недосяжною.

Особливо важливо, що суб’єкт (S) не лише інтегрує світ, а й підтримує власну інваріантність. Щоб залишатися собою, суб’єкт (S) змушений згортати безліч відмінностей у відносно стійку форму. Якщо ця форма втрачає стійкість, руйнується не лише цілісність самого суб’єкта (S), а й можливість бачити світ як узгоджену систему інваріантних структур ($\mathcal{I}$). У межах запропонованої рамки психічні захворювання можна інтерпретувати саме як частковий або повний розпад тих механізмів згортання й узгодження, які підтримують типову інваріантність \emph{Homo sapiens}.

У кімнаті з однорідним туманом це проявляється особливо ясно. Коли майже всі зовнішні інваріантні структури ($\mathcal{I}$) стають недоступними, суб’єкт (S) не перестає існувати повністю. Він продовжує утримувати мінімальну форму самості через внутрішні інваріантні структури ($\mathcal{I}$) — серцебиття, дихання, тілесне відчуття. Саме вони підтримують найпростішу інваріантність (S), навіть тоді, коли зовнішній світ майже зникає.

Водночас цей умоглядний експеримен не слід розуміти як нейтральний або нескінченно стабільний стан. Наукові ксперименти із сенсорною та перцептивною депривацією показують, що тривала втрата доступних зовнішніх інваріантних структур ($\mathcal{I}$) має руйнівний вплив на узгоджений суб’єкт (S). У дослідженнях, проведених у 1950-х роках під керівництвом Donald O. Hebb в McGill University, добровольці перебували в умовах майже повної сенсорної депривації: у невеликих кімнатах, із закритими очима, приглушеним звуком і мінімальним тактильним контактом. Уже через кілька годин у багатьох учасників виникали труднощі з послідовним мисленням, а приблизно через дві доби більшість повідомляла про галюцинації, втрату відчуття реальності та розпад звичних когнітивних структур. Через сильний психологічний вплив експерименти часто доводилося припиняти раніше запланованого часу; майже ніхто з учасників не міг витримати більше кількох діб~\cite{Hebb1958}.

У межах запропонованої рамки це можна інтерпретувати як наслідок того, що внутрішніх інваріантних структур ($\mathcal{I}$) виявляється недостатньо для довготривалого підтримання узгодженого суб’єкта (S). Коли зовнішні інваріантні структури ($\mathcal{I}$) зникають, (S) ще певний час зберігає себе через серцебиття, дихання і тілесну безперервність. Але якщо нові зовнішні інваріантні структури ($\mathcal{I}$) не з’являються, процес згортання починає втрачати стійкість. Узгоджений світ руйнується, а разом із ним починає руйнуватися і сам суб’єкт (S) як структура, здатна інтегрувати досвід.

Можливо, саме тому в умовах тривалої депривації виникають галюцинації. У цій перспективі їх можна розуміти не як випадкову помилку мозку, а як спробу суб’єкта (S) відновити відсутні інваріантні структури ($\mathcal{I}$) і знову побудувати хоча б мінімально узгоджений світ.

Це дозволяє стверджувати, що (S), як особливий інваріант, є результатом не одиничного акту, а як результат багаторівневого згортання, у якому дедалі складніші інваріантні структури ($\mathcal{I}$) інтегруються в більш стійку цілісність. На нижчих рівнях можуть існувати лише окремі інваріантні структури (I). На вищих рівнях з’являється можливість не лише утримувати ці структури, а й співвідносити їх між собою і зберігати безперервність самого процесу співвіднесення. Саме така структура і є суб’єктом (S).

У цьому сенсі (S) не є повним образом світу і не є незалежним від нього. Суб’єкт (S) і доступний світ виникають разом як дві сторони одного й того самого процесу згортання. Світ є системою інваріантних структур ($\mathcal{I}$), доступних для (S), а (S) є інваріантною структурою (I), яка здатна підтримувати цю систему як узгоджену.

\section{Межі гіпотези і відкриті питання}

Хоча автор достатньо обережно ставиться до сформульованих на цьому етапі рамок майбутньої теорії, уже тепер, спираючись на наведені міркування, можна спробувати висувати нові гіпотези та окреслювати перспективні напрями досліджень у таких галузях, як фундаментальна фізика, математика, кібернетика, нейробіологія і біологія загалом, геронтологія, пошук розуму (не лише позаземного), а також розвиток штучного інтелекту. Саме з цією метою і публікується ця робота.

Водночас автор визнає, що запропонована рамка є лише початковою спробою формалізувати інтуїтивні міркування про природу світу, суб’єкта \(S\) і їхню взаємодію. Вона не претендує ані на повноту, ані на статус остаточної істини. Навпаки, у її межах залишається багато відкритих питань, потенційних слабких місць і концептуальних напружень, які потребують подальшого аналізу.

Високий рівень абстракції та узагальнення, на якому побудована ця рамка, може виявитися як її силою, так і її слабкістю. З одного боку, він дозволяє охопити широкий спектр явищ і концепцій та потенційно стати підґрунтям для об’єднання під спільною теоретичною перспективою різних наукових напрямів. З іншого боку, саме цей рівень абстракції може ускладнити її перевірку, формалізацію і, відповідно, прийняття науковою спільнотою.

Одним із головних відкритих питань залишається статус самого поля відмінностей \(D\). Поки що воно вводиться як граничне концептуальне поняття, необхідне для того, щоб не ототожнювати реальність із уже сформованим світом об’єктів. Проте ще належить з’ясувати, чи може \(D\) бути строго визначене як фундаментальна онтологічна основа, чи воно має залишатися лише евристичним інструментом опису.

Не менш важливим є питання про природу інваріантних структур \(\mathcal{I}\). У статті вони описуються як умова побудови доступного світу для суб’єкта \(S\), однак поки не ясно, за якими саме критеріями певна структура може вважатися інваріантною, якою є межа між окремим інваріантом \(I\) та системою інваріантів \(\mathcal{I}\), і яким чином ці структури повинні бути формалізовані в межах майбутньої теорії.

Окремої уваги потребує і поняття згортання. У запропонованій рамці саме згортання пояснює перехід від поля відмінностей \(D\) до інваріантних структур \(\mathcal{I}\), а від них — до складніших рівнів організації, включно із суб’єктом \(S\). Проте механізм цього переходу поки що описано лише концептуально. Потрібно окремо з’ясувати, які саме формальні властивості має мати оператор згортання \(C\), за яких умов згортання зберігає стійкість, і чому одні лінії згортання можуть приводити до появи узгодженого суб’єкта \(S\), тоді як інші — ні.

Так само відкритим залишається питання про те, яким чином у межах цієї рамки слід розуміти простір і час. У статті вони інтерпретуються як похідні режими узгодження інваріантних структур \(\mathcal{I}\), але ця ідея поки що не має достатнього математичного апарату. Майбутня робота повинна показати, чи можливо послідовно вивести з такої перспективи хоча б деякі відомі властивості простору, часу і причинності.

Ще одним важливим питанням є статус самого суб’єкта \(S\). У запропонованій рамці він розглядається не як первинна сутність, а як інваріантна структура вищого порядку, що підтримує власну безперервність і узгоджує доступний світ. Однак поки що не ясно, де проходить межа між складною системою інваріантів \(\mathcal{I}\) і суб’єктом \(S\), за яких умов така система набуває самості, і яким чином цю межу можна було б описати у формально строгий спосіб.

Нарешті, запропонована рамка неминуче ставить питання про співвідношення між її концептуальним потенціалом і вимогами сучасної науки до перевірюваності. Якщо ця модель претендує на подальший розвиток, вона повинна буде породжувати або принаймні надихати конкретні дослідницькі програми, здатні давати нові інтерпретації вже відомих явищ або формулювати нові перевірювані гіпотези.

Саме тому цю статтю слід розглядати не як завершену теорію, а як перший крок: спробу окреслити нову концептуальну перспективу і ввести мову, яка дозволяє поставити питання, що виходять за межі традиційного поділу на об’єкт, суб’єкт, простір, час та інформацію.

\begin{quote}
Запропонована рамка є не відповіддю, а запрошенням до критичної дискусії про те, чи можуть світ, простір, час і суб’єкт \(S\) бути не первинними сутностями, а результатом багаторівневого згортання поля відмінностей \(D\).
\end{quote}

\section{Висновок}

Розглянутий у статті thought experiment із кімнатою, заповненою однорідним туманом, показує, що фізична наявність середовища ще не гарантує існування світу для суб’єкта \(S\). Навіть у середовищі, насиченому молекулами, фотонами та потенційними конфігураціями, світ може майже зникнути, якщо суб’єкт \(S\) не має доступу до інваріантних структур \(\mathcal{I}\), достатніх для побудови узгодженої картини.

На цій основі в статті було запропоновано розрізняти дані, інформацію і доступний світ. Дані охоплюють усе, що фізично присутнє. Інформація виникає лише там, де частина даних стає доступною у формі інваріантних структур \(\mathcal{I}\). Доступний світ є ще вужчим: це система інваріантних структур \(\mathcal{I}\), інтегрованих у стійку картину для конкретного суб’єкта \(S\).

Для опису більш фундаментального рівня було введено поняття поля відмінностей \(D\), яке не є ані готовим світом, ані простором, ані часом. Простір і час, у межах запропонованої рамки, постають як похідні режими узгодження інваріантних структур \(\mathcal{I}\). Саме тому вони можуть радикально змінюватися або майже зникати, коли відповідні інваріантні структури стають недоступними.

Далі була висунута гіпотеза багаторівневого згортання,
\[
D \to D' \to D'' \to \dots
\]
у межах якої кожен новий рівень зменшує кількість доступних відмінностей, але підвищує інтегрованість інваріантних структур \(\mathcal{I}\). У цій перспективі суб’єкт \(S\) не є первинною сутністю, а постає як інваріантна структура вищого порядку, що здатна інтегрувати доступні інваріантні структури \(\mathcal{I}\) в узгоджений світ і підтримувати власну безперервність.

Запропонована модель дозволяє по-новому поглянути на проблему суб’єктності, визначаючи її не як замкнену або самодостатню сутність, а як динамічний процес підтримки інваріантності в полі відмінностей \(D\). У цій перспективі суб’єкт \(S\) і світ не передують одне одному, а існують разом як дві сторони одного процесу: світ є системою інваріантних структур \(\mathcal{I}\), доступних для \(S\), а суб’єкт \(S\) є структурою, що інтегрує, утримує та узгоджує ці інваріанти.

Хоча автор достатньо обережно ставиться до сформульованих на цьому етапі рамок майбутньої теорії, уже зараз, спираючись на наведені міркування, можна пробувати висувати нові гіпотези та окреслювати перспективні напрями досліджень у таких галузях, як фундаментальна фізика, математика, кібернетика, нейробіологія і біологія загалом, геронтологія, пошук розуму (не лише позаземного), а також розвиток штучного інтелекту.

Таким чином, головний результат цієї роботи полягає не у побудові завершеної теорії, а у формулюванні нової концептуальної перспективи. Вона дозволяє розглядати суб’єкта \(S\), світ, простір, час і інформацію не як незалежні та первинні сутності, а як взаємопов’язані прояви одного процесу — процесу виникнення, згортання та підтримки інваріантності в полі відмінностей \(D\).

Саме з цією метою і публікується ця робота: не як завершена система, а як спроба окреслити мову і рамку, здатні стати основою для подальшої критичної, формальної та міждисциплінарної дискусії.

\bibliographystyle{apsrev4-2}
\begin{thebibliography}{9}

    \bibitem{Metzinger2003}
    T.~Metzinger,
    \textit{Being No One}
    (MIT Press, Cambridge, 2003).

    \bibitem{Bateson1972}
    G.~Bateson,
    \textit{Steps to an Ecology of Mind}
    (University of Chicago Press, Chicago, 1972).

    \bibitem{Friston2010}
    K.~Friston,
    ``The Free-Energy Principle: A Unified Brain Theory?''
    Nature Reviews Neuroscience \textbf{11}, 127--138 (2010).

    \bibitem{VarelaThompsonRosch1991}
    F.~Varela, E.~Thompson, and E.~Rosch,
    \textit{The Embodied Mind}
    (MIT Press, Cambridge, 1991).

    \bibitem{Mandelbrot1982}
    B.~Mandelbrot,
    \textit{The Fractal Geometry of Nature}
    (W. H. Freeman and Company, New York, 1982).

    \bibitem{Mandelbrot1977}
    B.~Mandelbrot,
    \textit{Fractals: Form, Chance and Dimension}
    (W. H. Freeman and Company, San Francisco, 1977).

    \bibitem{Hebb1958}
    D.~O.~Hebb,
    \textit{The Effects of Perceptual Isolation},
    Scientific American, 1958.
\end{thebibliography}

\vfill % Це змістить блок до самого низу сторінки
\begin{center}
    \rule{0.5\textwidth}{.4pt} \\ % Коротка горизонтальна лінія
    \vspace{2mm}
    \footnotesize
    © 2026 Віталій Дурнєв (Vitalii Durniev). Поширюється на умовах ліцензії CC BY-NC-ND 4.0. \\
    Цитування обов’язкове. Заборонено комерційне використання та поширення змінених версій без письмового дозволу автора.
\end{center}
\end{document}